% 引力波天文学（综述）
% license CCBYSA3
% type Wiki

本文根据 CC-BY-SA 协议转载翻译自维基百科\href{https://en.wikipedia.org/wiki/Gravitational-wave_astronomy}{相关文章}。


引力波天文学是天文学的一个分支，主要研究由天体物理源发出的引力波的探测与分析。\(^\text{[1]}\)

引力波是时空中极其微小的扭曲或涟漪，由大质量天体的加速运动所引起。它们产生于灾变性事件之中，例如双黑洞并合、双中子星合并、超新星爆炸，以及大爆炸后早期宇宙中的物理过程。研究引力波为观测宇宙提供了一种全新的方式，使人们能够深入了解物质在极端条件下的行为。类似于电磁辐射（如光波、无线电波、红外辐射和 X 射线）通过电磁场涨落传播能量的方式，引力辐射则是由相对较弱的引力场涨落所产生的能量传播。引力波的存在最早由奥利弗·赫维赛德（Oliver~Heaviside）于 1893 年提出，随后在 1905 年由昂利·庞加莱（Henri~Poincaré）作为电磁波的引力学对应物加以推测，最终在 1916 年被阿尔伯特·爱因斯坦（Albert~Einstein）根据广义相对论所预言。

1978 年，罗素·阿伦·赫尔斯（Russell~Alan~Hulse）与约瑟夫·胡顿·泰勒（Joseph~Hooton~Taylor~Jr.）通过观测两颗互相绕转的中子星，首次提供了引力波存在的实验性证据。二人因此工作获得了 1993 年诺贝尔物理学奖。2015 年，距爱因斯坦的预言近一个世纪后，人类首次直接探测到了由双黑洞并合产生的引力波信号，从而确认了这一长期难以捉摸的现象的存在，并开启了天文学的新纪元。此后，科学家又陆续探测到多起包括双黑洞并合、中子星碰撞及其他剧烈宇宙事件在内的引力波信号。如今，引力波的探测主要依靠激光干涉测量技术（laser interferometry），该技术通过测量两条相互垂直臂长随波动变化的微小差异来捕捉引力波。诸如 LIGO（Laser~Interferometer~Gravitational-wave~Observatory，激光干涉引力波天文台）、Virgo 与 KAGRA（Kamioka~Gravitational~Wave~Detector，神冈引力波探测器）等观测站都采用此项技术来探测来自遥远宇宙事件的微弱信号。LIGO 的联合创始人巴里·C·巴里什（Barry~C.~Barish）、基普·S·索恩（Kip~S.~Thorne）与赖纳·魏斯（Rainer~Weiss）因其在引力波天文学领域的开创性贡献而共同获得了 2017 年诺贝尔物理学奖。
